% #############################################################################
% Abstract Text
% !TEX root = ../main.tex
% #############################################################################
% use \noindent in first paragraph


\noindent Solar-powered vehicles, particularly boats, need effective energy conversion in rapidly changing conditions. The \ac{mppt} charger plays a critical role in energy conversion, ensuring the panels operate at their \ac{mpp}. This work focuses on the design, simulation and testing of a fast response \ac{mppt} charger for \ac{tsb}.

The prototype was designed to charge a lithium-ion battery of 43.2~V nominal voltage. The proposed charger is based on a synchronous boost converter operating at 100 kHz controlled by a \ac{mcu}. 

Different \ac{mppt} techniques were analyzed to identify a suitable solution for the dynamic operating conditions of the vessel. Balancing tracking speed and accuracy led to the implementation of a \ac{peo} algorithm. The converter and \ac{pv} system were simulated in both LTspice and Matlab to validate the component selection and algorithm parameters. A custom \ac{pcb} was designed, manufactured and soldered by hand. 

Experimental results show that the implemented \ac{mppt} algorithm can respond rapidly to changes in irradiance or temperature, with the measured tracking response being suitable for the environmental variations identified. The prototype achieved a measured efficiency of up to 96\%. Beyond performance, a computer application was also developed for configuration and data analysis. The resulting system provides a dedicated, configurable, and low-cost \ac{mppt} charger for \ac{tsb} and a platform for future improvements.

\vspace{1cm}